% Source: uncertain_and_unreliable_evidence.tex
\begin{notechapter}
\chapter{Uncertain and Unreliable Evidence}\label{ch:uncertain-evidence}

\section{Introduction}
In Bayesian reasoning, evidence is used to update beliefs about
the state of the world. However, evidence can be \textbf{uncertain}
(noisy or imprecise) or \textbf{unreliable} (potentially incorrect).
This document explains these concepts and provides fully developed
examples.

\section{Uncertain Evidence}
Uncertain evidence refers to observations that are not perfectly
known. For example, a sensor reading might be noisy. We model
uncertain evidence using probability distributions.

\subsection{Bayesian Framework}
Given a true state \( x \) and uncertain evidence \( y \), we use
Bayes' rule to update our belief:
\[
p(x | y) = \frac{p(y | x) \, p(x)}{p(y)}
\]
Here:
\begin{itemize}
    \item \( p(x) \): Prior belief about \( x \).
    \item \( p(y | x) \): Likelihood of observing \( y \) given \( x \).
    \item \( p(x | y) \): Posterior belief about \( x \) after observing \( y \).
\end{itemize}

\section{Unreliable Evidence and Virtual Evidence}
Unreliable evidence refers to observations that might be incorrect
or misleading. For example, a witness might lie. We model
unreliability by introducing a \textbf{virtual evidence term} that encodes
the decision-maker's subjective interpretation of the evidence.

\subsection{Virtual Evidence (Barber's Framework)}
Virtual evidence is used to handle unreliable evidence by replacing
the original likelihood term \( p(G | A) \) with a virtual evidence
term \( p(H | A) \). This term encodes the decision-maker's belief
about the true state \( A \) based on their interpretation of the
unreliable evidence.

For example, if the decision-maker believes there is an 80\% chance
the alarm sounded (\( A = \text{tr} \)) and a 20\% chance it did not
(\( A = \text{fa} \)), the virtual evidence term is:
\[
p(H | A) = 
\begin{cases}
0.8 & \text{if } A = \text{tr}, \\
0.2 & \text{if } A = \text{fa}.
\end{cases}
\]

\subsection{Effect on the Joint Distribution}
The virtual evidence term \( p(H | A) \) is incorporated into the
joint distribution as a factor:
\[
p(B, A, H, W) = p(A | B) \, p(B) \, p(W | A) \, p(H | A).
\]
This modifies the joint distribution to reflect the decision-maker's
interpretation of the unreliable evidence.

\subsection{Posterior Update}
When computing the posterior \( p(B = \text{tr} | W = \text{tr}, H) \),
the virtual evidence term \( p(H | A) \) acts as a likelihood that
favors certain states of \( A \) over others. For example, in the
case above, \( p(H | A) \) favors \( A = \text{tr} \) by a factor
of 4 compared to \( A = \text{fa} \).

\section{Example: Holmes and the Alarm}
Suppose Holmes receives unreliable evidence from Mrs Gibbon about
whether an alarm sounded (\( A \)). He believes there is an 80\%
chance the alarm sounded and a 20\% chance it did not. We model
this using a virtual evidence term \( p(H | A) \):
\[
p(H | A) = 
\begin{cases}
0.8 & \text{if } A = \text{tr}, \\
0.2 & \text{if } A = \text{fa}.
\end{cases}
\]

\subsection{Step 1: Define the Joint Distribution}
The joint distribution is:
\[
p(B, A, H, W) = p(A | B) \, p(B) \, p(W | A) \, p(H | A).
\]

\subsection{Step 2: Compute the Posterior}
Using Bayes' rule, and since the alarm state \( A \) is not observed, we must marginalize over it:
\[
p(B = \text{tr} | W = \text{tr}, H) \propto p(B = \text{tr}) \sum_{a \in \{\text{tr},\text{fa}\}} p(A = a | B = \text{tr}) \, p(W = \text{tr} | A = a) \, p(H | A = a),
\]
and similarly for \( B = \text{fa} \); normalizing the two values gives the posterior. The evidence
\( W \) and the virtual evidence \( H \) only reach \( B \) through the unobserved alarm \( A \).

\section{Example 1: Is It Raining?}
Suppose we want to determine if it's raining (\( x = 1 \)) or not
(\( x = 0 \)). We have two pieces of evidence:
\begin{itemize}
    \item A weather app says it's raining (\( y_1 = 1 \)), but it's
          correct only 80\% of the time.
        \item A friend says it's not raining (\( y_2 = 0 \)), but he
          is not always reliable: with probability 0.3 he is in an
          ``unreliable'' state, in which he lies 30\% of the time (so
          overall he lies \( 0.3 \times 0.3 = 9\% \) of the time).
\end{itemize}

\subsection{Step 1: Model Uncertain Evidence}
The weather app provides uncertain evidence. We model its
likelihood:
\[
p(y_1 = 1 | x = 1) = 0.8, \quad p(y_1 = 1 | x = 0) = 0.2
\]

\subsection{Step 2: Model Unreliable Evidence (Latent Reliability Variable)}
The friend provides unreliable evidence. Instead of a virtual evidence term, we
model the unreliability explicitly with a latent reliability
variable \( r \) that is marginalized out, where \( r = 1 \) means the friend is reliable,
and \( r = 0 \) means he is unreliable.

\begin{itemize}
    \item \textbf{If the friend is reliable (\( r = 1 \)):}
    \begin{itemize}
        \item He will tell the truth.
        \item If it's raining (\( x = 1 \)), he will say it's raining (\( y_2 = 1 \)).
        \item If it's not raining (\( x = 0 \)), he will say it's not raining (\( y_2 = 0 \)).
    \end{itemize}
    So:
    \[
    p(y_2 = 1 | x = 1, r = 1) = 1, \quad p(y_2 = 0 | x = 1, r = 1) = 0
    \]
    \[
    p(y_2 = 0 | x = 0, r = 1) = 1, \quad p(y_2 = 1 | x = 0, r = 1) = 0
    \]

    \item \textbf{If the friend is unreliable (\( r = 0 \)):}
    \begin{itemize}
        \item He might lie.
        \item Suppose he lies 30\% of the time.
    \end{itemize}
    So:
    \[
    p(y_2 = 1 | x = 1, r = 0) = 0.7, \quad p(y_2 = 0 | x = 1, r = 0) = 0.3
    \]
    \[
    p(y_2 = 0 | x = 0, r = 0) = 0.7, \quad p(y_2 = 1 | x = 0, r = 0) = 0.3
    \]
\end{itemize}

\subsection{Step 3: Combine Evidence}
Assume a prior \( p(x = 1) = 0.5 \) and \( p(r = 1) = 0.7 \). Using
Bayes' rule, we compute the posterior:
\[
p(x = 1 | y_1 = 1, y_2 = 0) \propto p(y_1 = 1 | x = 1) \,
p(y_2 = 0 | x = 1) \, p(x = 1)
\]
\[
p(x = 0 | y_1 = 1, y_2 = 0) \propto p(y_1 = 1 | x = 0) \,
p(y_2 = 0 | x = 0) \, p(x = 0)
\]

First, compute \( p(y_2 = 0 | x = 1) \) and \( p(y_2 = 0 | x = 0) \):
\[
p(y_2 = 0 | x = 1) = p(y_2 = 0 | x = 1, r = 1) p(r = 1) +
p(y_2 = 0 | x = 1, r = 0) p(r = 0)
\]
\[
p(y_2 = 0 | x = 1) = 0 \times 0.7 + 0.3 \times 0.3 = 0.09
\]
\[
p(y_2 = 0 | x = 0) = p(y_2 = 0 | x = 0, r = 1) p(r = 1) +
p(y_2 = 0 | x = 0, r = 0) p(r = 0)
\]
\[
p(y_2 = 0 | x = 0) = 1 \times 0.7 + 0.7 \times 0.3 = 0.91
\]

Now, compute the unnormalized posteriors:
\[
p(x = 1 | y_1 = 1, y_2 = 0) \propto 0.8 \times 0.09 \times 0.5 = 0.036
\]
\[
p(x = 0 | y_1 = 1, y_2 = 0) \propto 0.2 \times 0.91 \times 0.5 = 0.091
\]

Normalize to obtain the final probabilities:
\[
p(x = 1 | y_1 = 1, y_2 = 0) = \frac{0.036}{0.036 + 0.091} \approx 0.283
\]
\[
p(x = 0 | y_1 = 1, y_2 = 0) = \frac{0.091}{0.036 + 0.091} \approx 0.717
\]

\section{Example 2: Medical Diagnosis}
Suppose a patient is tested for a disease (\( x = 1 \) if diseased,
\( x = 0 \) otherwise). The test result (\( y \)) is positive, but
the test has a 5\% false positive rate and a 10\% false negative rate.

\subsection{Step 1: Model Uncertain Evidence}
The test provides uncertain evidence. We model its likelihood:
\[
p(y = 1 | x = 1) = 0.9, \quad p(y = 1 | x = 0) = 0.05
\]

\subsection{Step 2: Incorporate Prior}
Assume the disease prevalence is 1\%, so \( p(x = 1) = 0.01 \).

\subsection{Step 3: Compute Posterior}
Using Bayes' rule:
\[
p(x = 1 | y = 1) = \frac{p(y = 1 | x = 1) \, p(x = 1)}{p(y = 1)}
\]
\[
p(y = 1) = p(y = 1 | x = 1) \, p(x = 1) + p(y = 1 | x = 0) \, p(x = 0)
\]
\[
p(y = 1) = 0.9 \times 0.01 + 0.05 \times 0.99 = 0.0585
\]
\[
p(x = 1 | y = 1) = \frac{0.9 \times 0.01}{0.0585} \approx 0.153
\]
The probability of having the disease given a positive test is 15.3\%.

\section{Conclusion}
Bayesian reasoning provides a principled way to handle uncertain
and unreliable evidence. By modeling uncertainty and unreliability,
we can update our beliefs in a rigorous and interpretable manner.
\end{notechapter}
